\chapter{EXPECTED OUTCOME}

The primary goal of this project is to develop and validate a research-grade prototype for explainable, wind-aware air quality forecasting. The proposed system serves as a decision-support tool by bridging the gap between raw environmental data and actionable policy insights.

\section{Technical Outcomes}

The expected technical outcomes of the project are as follows:

\begin{itemize}
    \item \textbf{Robust Data Pipeline:}  
    A fully automated preprocessing framework for multi-source data ingestion, timestamp synchronization, outlier detection, missing value handling, and normalization to support high-quality model training.

    \item \textbf{Wind-Aware Dynamic Graph Engine:}  
    A dynamic graph construction method that updates station-to-station relationships by integrating spatial distance, pollutant correlations, and wind direction and speed information.

    \item \textbf{Hybrid Forecasting Framework:}  
    A forecasting framework consisting of traditional baseline models such as LSTM and GRU, along with advanced spatio-temporal graph neural network models such as GCN-GRU and GAT-GRU for performance benchmarking.

    \item \textbf{Explainability and Decision Support Module:}  
    A module that interprets model outputs to provide human-readable insights, such as important meteorological features, influential neighboring stations, and temporal factors contributing to high-pollution events.

    \item \textbf{Interactive Analytics Dashboard:}  
    A centralized visualization platform for displaying air quality forecasts, spatial risk maps, historical trends, and interpretability reports for end-users.
\end{itemize}

